% Program: PortfolioCh4S1-2.tex
% Author: Carmen Beltran, Brendan Butler, Teddy Wachtler
%
% Description: Portfolio Assignment 7: Sections 4.1 and 4.2
%%%%%%%%%%%%%%%%%%%%%%%%%%%%%%%%%%%%%%%%%%%%%%%%%%%%%%%%%%%

\documentclass{article}
\usepackage[dvips]{graphicx}
\usepackage{a4wide}
\usepackage{amsmath}
\usepackage{euscript}
\usepackage{amssymb}
\usepackage{amsthm}
\usepackage{amsopn}

\theoremstyle{definition}
\newtheorem*{definition}{Definition}
\newtheorem{theorem}{Theorem}

\renewcommand{\qed}{\blacksquare}
\newcommand{\powset}{{\cal P}}
\newcommand{\img}{_{*}}
\newcommand{\pimg}{^{*}}

\newcommand{\vv}{\ensuremath{\vec{v}}}
\newcommand{\vu}{\ensuremath{\vec{u}}}
\newcommand{\vw}{\ensuremath{\vec{w}}}
\newcommand{\vx}{\ensuremath{\vec{x}}}
\newcommand{\vy}{\ensuremath{\vec{y}}}
\newcommand{\vb}{\ensuremath{\vec{b}}}
\newcommand{\vo}{\ensuremath{\vec{0}}}
\newcommand{\va}{\ensuremath{\vec{a}}}
\newcommand{\ve}{\ensuremath{\vec{e}}}

\newcommand{\R}{\mathbb{R}}
\newcommand{\Z}{\mathbb{Z}}
\newcommand{\C}{\mathbb{C}}
\newcommand{\N}{\mathbb{N}}
\newcommand{\Q}{\mathbb{Q}}

%%%%%%%%%%%%%%%%%%%%%%%%%%%%%%%%%%%%%%%%%%%%%%%%%%%%%%%%%%%

\begin{document}
	\begin{center}
		\Large{Math 251W: Foundations of Advanced Mathematics}

		\normalsize Portfolio problems from section 3.4

		\hfill Name: Carmen Beltran, Brendan Butler, Teddy Wachtler
		\hrule
		\vspace{0.4cm}
		\vspace{0.75cm}
	\end{center}

%%%%%%%%%%%%%%%%%%%%%%%%%%%%%%%%%%%%%%%%%%%%%%%%%%%%%%%%%%%

	\vspace{1cm}
	\noindent Problem 4.1.8
	\vspace{0.3cm}

	\underline{proposition:}
	$A, B\subseteq X, \chi_A = \chi_B$ iff $A = B$.

	\vspace{.2cm}
	\fbox{proof} (Direct)
	\vspace{.15cm}

	Let $A, B$, and $X$ be arbitrary sets so that $A, B \subseteq X$.
	Let $\chi_A : X \rightarrow \{0, 1\}$ be the characteristic function of the subset $A$ in $X$.
	Let $\chi_B : X \rightarrow \{0, 1\}$ be the characteristic function of the subset $B$ in $X$.
	We ask whether $\chi_A = \chi_B$ if and only if $A = B$.

	Consider the case where $\chi_A = \chi_B$.
	Let $a$ be an arbitrary element in $A$.
	By the definition of the characteristic function, $\chi_A(a) = 1$.
	Since $\chi_A = \chi_B$, we see that $\chi_B(a) = 1$.
	By the definition of the characteristic function, we see that $a \in B$.
	By the definition of subsets, since $a$ was an arbitrary element of $A$, we see that $A \subseteq B$.
	Let $b$ be an arbitrary element in $B$.
	By symmetry, we see $\chi_B(b) = 1$, $\chi_A(b) = 1$, $b \in A$ and $B \subseteq A$.
	Therefore, by the definition of set equality, since $A \subseteq B$ and $B \subseteq A$, we see that if $\chi_A = \chi_B$, then $A = B$.

	Consider the case where $A = B$.
	Let $x$ be an arbitrary element in $X$.
	By the definition of the characteristic function, we see that $\chi_A(x) = 1$ if and only if $x \in A$, and, $\chi_A(x) = 0$ if and only if $x \not \in A$.
	Similarly, we see that $\chi_B(x) = 1$ if and only if $x \in B$.
	By the definition of set equality, we see that $x \in A$ if and only if $x \in B$.
	Therefore, we see that $\chi_A(x) = \chi_B(x)$.
	Therefore, since $x$ was an arbitrary element of $X$, we see that if $A = B$ then $\chi_A = \chi_B$.

	Therefore, since if $\chi_A = \chi_B$ then $A = B$, and if $A = B$ then $\chi_A = \chi_B$, we see $\chi_A = \chi_B$ if and only if $A = B$.
	Therefore, for all sets $A, B$, and $X$, if $A, B \subseteq X$, then $\chi_A = \chi_B$ if and only if $A = B$.
	$\qed$.

	\vspace{1cm}

%%%%%%%%%%%%%%%%%%%%%%%%%%%%%%%%%%%%%%%%%%%%%%%%%%%%%%%%%%%

	\noindent Problem 4.2.11
	\vspace{0.3cm}

	\underline{proposition:}
	Given $f : A \rightarrow B$ is a function, and $P, Q \subseteq A, f \img(P) - f \img(Q) \subseteq f \img(P - Q)$.

	\vspace{.2cm}
	\fbox{proof} (Direct)
	\vspace{.15cm}

	Let $A, B, P$ and $Q$ be arbitrary sets so that $P, Q \subseteq A$.
	Let $f : A \rightarrow B$ be a function.
	Let $x$ be an arbitrary element of $f \img(P) - f \img(Q)$.
	By the definition of difference of sets, we see $x \in f \img(P)$ and $x \not \in f \img(Q)$.
	By the definition of set images, since $x \in f \img(P)$, there exists a particular element $y \in P$ so that $f(y) = x$.
	By the definition of set images, since $x \not \in f \img(Q)$, we see that for all elements $z \in Q$, $f(z) \neq x$.
	Therefore, since $f(y) = x$, we see $y \not \in Q$.
	By the definition of difference of sets, since $y \in P$ and $y \not \in Q$, we see $y \in P - Q$.
	By the definition of set images, since $y \in P - Q$, we see $f(y) \in f \img(P - Q)$.
	Since $f(y) = x$, we see $x \in f \img(P - Q)$.
	Therefore, by the properties of subsets, since $x$ was an arbitrary element of $f \img(P) - f \img(Q)$, we see that $f \img(P) - f \img(Q) \subseteq f \img(P - Q)$.
	Therefore, for all sets $A, B, P$ and $Q$, if $f : A \rightarrow B$ is a function and $P, Q \subseteq A$, then $f \img(P) - f \img(Q) \subseteq f \img(P - Q)$.
	$\qed$.

	\vspace{1cm}

%%%%%%%%%%%%%%%%%%%%%%%%%%%%%%%%%%%%%%%%%%%%%%%%%%%%%%%%%%%

\end{document}